% GENERATED FILE. Edit canonical lessons, then run npm run generate:books.
% canonical-source-hash: fnv1a64:145cc13043abbd61
% canonical-lessons: MR-C522-divankhana, MR-C522-sayangrh, MR-C522-gacci, MR-C522-pharsi, MR-C522-tangne

\chapter{Rooms of the House}
\label{ch:mr-rooms-of-the-house}
\hlchaptermodality{522}

\begin{chapteropening}
\textbf{By the end of this chapter:} \emph{I can name the rooms of a house and hang things on the walls.}

Complete the last lesson of chapter 522: I can name the rooms of a house and hang things on the walls.
\end{chapteropening}

\section[\texorpdfstring{\mr{दिवाणखाना} (divāṇkhānā)}{दिवाणखाना (divāṇkhānā)}]{\mr{दिवाणखाना} (divāṇkhānā) — a living room}

\label{lesson:MR-C522-divankhana}



\begin{quote}
Before the new one: say the Marathi for a saw, then the Marathi for to plaster.
\end{quote}

\subsection*{You'll want to know: \mr{दिवाणखाना}}
\textbf{\mr{दिवाणखाना}} — \emph{divāṇkhānā} — \textquotedblleft{}a living room\textquotedblright{}.

\begin{quote}
\mr{पाहुणे} \mr{दिवाणखान्यात} \mr{बसले} \mr{होते}.
\end{quote}

\begin{quote}
\emph{pāhuṇe divāṇkhānyāt basle hote.}
\end{quote}

\begin{quote}
The guests were sitting in the living room.
\end{quote}

Masculine.

The oblique is \mr{दिवाणखान्या}-: \mr{दिवाणखान्यात}.

\begin{cousinweb}
\textbf{\mr{दिवाणखाना}} (\emph{divāṇkhānā}) is a loan from Persian \emph{dīvān-khāna}, \textquotedblleft{}the hall where guests are received\textquotedblright{}, from \emph{dīvān}, \textquotedblleft{}court\textquotedblright{}, and \emph{khāna}, \textquotedblleft{}house\textquotedblright{}.
\end{cousinweb}

\begin{grammarlens}[title={What you've built}]
One more word for describing what you lived and felt.
\end{grammarlens}

\subsection*{Your turn}
\begin{itemize}
\raggedright
  \item \emph{Say it:} \emph{divāṇkhānā}
  \item \emph{Say it:} \emph{divāṇkhānā}, once more
  \item \emph{Say it:} \emph{pāhuṇe divāṇkhānyāt basle hote.} — The guests were sitting in the living room.
  \item \emph{Recall:} say the Marathi for a saw, then the Marathi for to plaster, then say \emph{divāṇkhānā} again
\end{itemize}

\begin{tcolorbox}[breakable,colback=teal!4,colframe=teal!35!black,title={Before you move on}]
What does \mr{दिवाणखाना} mean? (\textquotedblleft{}A living room\textquotedblright{}.) Say it once more.
\end{tcolorbox}
\section[\texorpdfstring{\mr{शयनगृह} (śayangṛh)}{शयनगृह (śayangṛh)}]{\mr{शयनगृह} (śayangṛh) — a bedroom}

\label{lesson:MR-C522-sayangrh}



\begin{quote}
Before the new one: say the Marathi for to plaster, then the Marathi for a living room.
\end{quote}

\subsection*{You'll want to know: \mr{शयनगृह}}
\textbf{\mr{शयनगृह}} — \emph{śayangṛh} — \textquotedblleft{}a bedroom\textquotedblright{}.

\begin{quote}
\mr{आमच्या} \mr{घरात} \mr{दोन} \mr{शयनगृहं} \mr{आहेत}.
\end{quote}

\begin{quote}
\emph{āmcyā gharāt don śayangṛhaṁ āhet.}
\end{quote}

\begin{quote}
There are two bedrooms in our house.
\end{quote}

Neuter. Plural: \mr{शयनगृहं}, \emph{śayangṛhaṁ}.

\begin{cousinweb}
\textbf{\mr{शयनगृह}} (\emph{śayangṛh}) is a tatsama compound of Sanskrit \emph{śayana}, \textquotedblleft{}lying down, sleeping\textquotedblright{}, and \emph{gṛha}, \textquotedblleft{}house\textquotedblright{} (Monier-Williams).
\end{cousinweb}

\begin{grammarlens}[title={What you've built}]
One more word for describing what you lived and felt.
\end{grammarlens}

\subsection*{Your turn}
\begin{itemize}
\raggedright
  \item \emph{Say it:} \emph{śayangṛh}
  \item \emph{Say it:} \emph{śayangṛh}, once more
  \item \emph{Say it:} \emph{āmcyā gharāt don śayangṛhaṁ āhet.} — There are two bedrooms in our house.
  \item \emph{Recall:} say the Marathi for to plaster, then the Marathi for a living room, then say \emph{śayangṛh} again
\end{itemize}

\begin{tcolorbox}[breakable,colback=teal!4,colframe=teal!35!black,title={Before you move on}]
What does \mr{शयनगृह} mean? (\textquotedblleft{}A bedroom\textquotedblright{}.) Say it once more.
\end{tcolorbox}
\section[\texorpdfstring{\mr{गच्ची} (gaccī)}{गच्ची (gaccī)}]{\mr{गच्ची} (gaccī) — a terrace; a flat roof}

\label{lesson:MR-C522-gacci}



\begin{quote}
Before the new one: say the Marathi for a living room, then the Marathi for a bedroom.
\end{quote}

\subsection*{You'll want to know: \mr{गच्ची}}
\textbf{\mr{गच्ची}} — \emph{gaccī} — \textquotedblleft{}a terrace; a flat roof\textquotedblright{}.

\begin{quote}
\mr{उन्हाळ्यात} \mr{आम्ही} \mr{गच्चीवर} \mr{झोपायचो}.
\end{quote}

\begin{quote}
\emph{unhāḷyāt āmhī gaccīvar jhopāyco.}
\end{quote}

\begin{quote}
In summer we used to sleep on the terrace.
\end{quote}

Feminine.

\begin{cousinweb}
\textbf{\mr{गच्ची}} (\emph{gaccī}) is the word for a flat roof or terrace; its origin is not settled.
\end{cousinweb}

\begin{grammarlens}[title={What you've built}]
One more word for describing what you lived and felt.
\end{grammarlens}

\subsection*{Your turn}
\begin{itemize}
\raggedright
  \item \emph{Say it:} \emph{gaccī}
  \item \emph{Say it:} \emph{gaccī}, once more
  \item \emph{Say it:} \emph{unhāḷyāt āmhī gaccīvar jhopāyco.} — In summer we used to sleep on the terrace.
  \item \emph{Recall:} say the Marathi for a living room, then the Marathi for a bedroom, then say \emph{gaccī} again
\end{itemize}

\begin{tcolorbox}[breakable,colback=teal!4,colframe=teal!35!black,title={Before you move on}]
What does \mr{गच्ची} mean? (\textquotedblleft{}A terrace; a flat roof\textquotedblright{}.) Say it once more.
\end{tcolorbox}
\section[\texorpdfstring{\mr{फरशी} (pharśī)}{फरशी (pharśī)}]{\mr{फरशी} (pharśī) — a floor tile; a stone floor}

\label{lesson:MR-C522-pharsi}



\begin{quote}
Before the new one: say the Marathi for a bedroom, then the Marathi for a terrace.
\end{quote}

\subsection*{You'll want to know: \mr{फरशी}}
\textbf{\mr{फरशी}} — \emph{pharśī} — \textquotedblleft{}a floor tile; a stone floor\textquotedblright{}.

\begin{quote}
\mr{नवीन} \mr{घरात} \mr{पांढरी} \mr{फरशी} \mr{आहे}.
\end{quote}

\begin{quote}
\emph{navīn gharāt pāṇḍhrī pharśī āhe.}
\end{quote}

\begin{quote}
The new house has a white tiled floor.
\end{quote}

Feminine.

\begin{cousinweb}
\textbf{\mr{फरशी}} (\emph{pharśī}) is usually traced to Persian \emph{farš}, \textquotedblleft{}carpet, floor\textquotedblright{}, which goes back to Arabic \emph{farš}.
\end{cousinweb}

\begin{grammarlens}[title={What you've built}]
One more word for describing what you lived and felt.
\end{grammarlens}

\subsection*{Your turn}
\begin{itemize}
\raggedright
  \item \emph{Say it:} \emph{pharśī}
  \item \emph{Say it:} \emph{pharśī}, once more
  \item \emph{Say it:} \emph{navīn gharāt pāṇḍhrī pharśī āhe.} — The new house has a white tiled floor.
  \item \emph{Recall:} say the Marathi for a bedroom, then the Marathi for a terrace, then say \emph{pharśī} again
\end{itemize}

\begin{tcolorbox}[breakable,colback=teal!4,colframe=teal!35!black,title={Before you move on}]
What does \mr{फरशी} mean? (\textquotedblleft{}A floor tile; a stone floor\textquotedblright{}.) Say it once more.
\end{tcolorbox}
\section[\texorpdfstring{\mr{टांगणे} (ṭāṅgṇe)}{टांगणे (ṭāṅgṇe)}]{\mr{टांगणे} (ṭāṅgṇe) — to hang up}

\label{lesson:MR-C522-tangne}



\begin{quote}
Before the new one: say the Marathi for a terrace, then the Marathi for a floor tile.
\end{quote}

\subsection*{You'll want to know: \mr{टांगणे}}
\textbf{\mr{टांगणे}} — \emph{ṭāṅgṇe} — \textquotedblleft{}to hang up\textquotedblright{}.

\begin{quote}
\mr{मी} \mr{दिवाणखान्याच्या} \mr{भिंतीवर} \mr{एक} \mr{चित्र} \mr{टांगलं}.
\end{quote}

\begin{quote}
\emph{mī divāṇkhānyācyā bhintīvar ek citra ṭāṅglaṁ.}
\end{quote}

\begin{quote}
I hung a picture on the living room wall.
\end{quote}

Past, with -\mr{ने} (he …): \mr{त्याने} \mr{चित्र} \mr{टांगलं}, \emph{tyāne citra ṭāṅglaṁ}.

The place takes -\mr{वर}: \mr{भिंतीवर} \mr{टांगणे}.

\begin{cousinweb}
\textbf{\mr{टांगणे}} (\emph{ṭāṅgṇe}) is inherited: CDIAL lists Marathi \emph{ṭā̃gṇẽ}, \textquotedblleft{}to hang up\textquotedblright{}, under \emph{*ṭaṅg}, \textquotedblleft{}hang up\textquotedblright{} (CDIAL 5436).
\end{cousinweb}

\begin{grammarlens}[title={What you've built}]
One more word for describing what you lived and felt.
\end{grammarlens}

\subsection*{Your turn}
\begin{itemize}
\raggedright
  \item \emph{Say it:} \emph{ṭāṅgṇe}
  \item \emph{Say it:} \emph{ṭāṅgṇe}, once more
  \item \emph{Say it:} \emph{mī divāṇkhānyācyā bhintīvar ek citra ṭāṅglaṁ.} — I hung a picture on the living room wall.
  \item \emph{Recall:} say the Marathi for a terrace, then the Marathi for a floor tile, then say \emph{ṭāṅgṇe} again
\end{itemize}

\begin{tcolorbox}[breakable,colback=teal!4,colframe=teal!35!black,title={Before you move on}]
What does \mr{टांगणे} mean? (\textquotedblleft{}To hang up\textquotedblright{}.) Say it once more.
\end{tcolorbox}
